\section{Introduction}

Human assessment is the slowest and most expensive part of a TREC track. Large
language models (LLMs) can grade answers for a fraction of the cost, and several
studies find that system rankings built from their grades often agree with rankings
built from human
grades~\cite{faggioli2023perspectives,thomas2024predict,upadhyay2024umbrela,pradeep2025greatnugget}.
Where such grades can stand in for human assessment, and where they break, is
still open. TREC 2026 AutoJudge turns that question into a shared
task~\cite{autojudge2026}. A participant builds a judge. The judge reads the
anonymized runs of a host track and ranks the systems that produced them. It is
scored by how closely its ranking matches the ranking that human assessors
produce.

We built two judges that start from opposite assumptions. \judge{NuggetRecall}
assumes that the request alone says what a good answer must contain. An LLM writes
an answer key of facts before it sees any response, in the tradition of nugget
evaluation~\cite{voorhees2003qa,lin2006methods,pradeep2024initial}, and then
checks each response against that key. \judge{Corroborate} uses no LLM. It
assumes that a claim many independent teams make is probably one that matters. It
pools the sentences of all responses to a topic, counts how many distinct teams
state each claim, and scores a response by the claims it shares with other teams.
The idea comes from pooling in retrieval evaluation, where the merged output of
many systems decides which documents get judged~\cite{zobel1998reliable}, and from
work that ranks systems from that pool without any human
judgments~\cite{soboroff2001ranking}.

This paper makes four contributions.
\begin{itemize}
  \item Two complete judges, each submitted as data and as code.
    \judge{Corroborate} never lets a team vote for its own claims. A claim needs
    two teams to enter the bank, and it is weighted by the number of teams other
    than the scored one that state it. This is a team-level version of the
    leave-one-run-out test of pool bias of~\citet{zobel1998reliable}.
  \item Evidence from the pilots that the two judges do best on different human
    measures. On long reports \judge{NuggetRecall} is our stronger judge on nugget
    coverage and \judge{Corroborate} on F1. On short answers \judge{Corroborate}
    trails on both truth measures.
  \item Measurements of what matters inside an LLM judge. The judge model mattered
    most. Ranking runs within each topic before averaging gave a small gain on
    short answers and a small loss on long reports.
  \item Six ways our judges failed, or would have failed, during development
    without raising an error. Several produced a leaderboard that looked plausible.
    We describe the checks that caught them.
\end{itemize}

Every pilot score in this paper comes from the organizers' evaluation on TIRA,
and every comparison with other entries is informal. Test-set results will be
added once the organizers release them.
